
\SGASubsection{1.2}{Global lifting of finite groups}
\SGARefTarget{sga3:xv:subsection:1:2}

\SGAStatement{\SGAInlineRefTarget{sga3:xv:lemma:1-4}Lemma 1.4}{Let \(A\) be a local ring, separated and
complete for the topology defined by its maximal ideal \(\mathfrak m\),
let
\[
  S=\operatorname{Spec}(A),\qquad
  S_n=\operatorname{Spec}(A/\mathfrak m^n).
\]
Then, for every prescheme \(X\) (respectively, for every \(S\)-prescheme
\(X\)), the canonical map
\[
  \tag{$\ast$}\label{sga3:xv:lemma:1-4:equation:ast}
  \operatorname{Hom}(S,X)
  \longrightarrow
  \varprojlim_n\operatorname{Hom}(S_n,X)
\]
(respectively,
\[
  \tag{$\ast\ast$}\label{sga3:xv:lemma:1-4:equation:ast-ast}
  \Gamma(X/S)
  \longrightarrow
  \varprojlim_n\Gamma(X_n/S_n),
  \qquad X_n=X\times_S S_n)
\]
is bijective.}
\SGARefTarget{sga3:xv:lemma:1:4}

Statement \((\ast\ast)\) is an easy consequence of \((\ast)\). Let us
prove \((\ast)\). Let
\[
  u_n:S_n\longrightarrow X\qquad(n\in\mathbb N)
\]
be a coherent system of morphisms. The image \(y\) of the closed point
of \(S_n\) is independent of \(n\), and \(u_n\) factors through
\(\operatorname{Spec}\mathcal O_y\). On passing to the inverse limit,
the morphisms \(u_n\) define a ring homomorphism
\[
  \widetilde u:\mathcal O_y
  \longrightarrow
  \varprojlim_n(A/\mathfrak m^n)=A.
\]
This proves that \((\ast)\) is surjective; it is injective as soon as
\(A\) is separated for the \(\mathfrak m\)-adic topology.

\SGAStatement{\SGAInlineRefTarget{sga3:xv:corollary:1-5}Corollary 1.5}{Let \(A\) be a complete noetherian local
ring, let \(\mathfrak m\) be its maximal ideal, let
\[
  S=\operatorname{Spec}(A),\qquad
  S_n=\operatorname{Spec}(A/\mathfrak m^n),
\]
let \(X\) be a finite scheme over \(S\), and let \(Y\) be an
\(S\)-prescheme. Then the canonical map
\[
  \operatorname{Hom}_S(X,Y)
  \longrightarrow
  \varprojlim_n\operatorname{Hom}_{S_n}(X_n,Y_n)
\]
(where \(X_n=X\times_S S_n\), and similarly for \(Y_n\)) is
bijective.}
\SGARefTarget{sga3:xv:corollary:1:5}

Indeed, it follows from EGA II~6.2.5 that \(X\) is a finite disjoint
sum of local \(S\)-schemes finite over \(S\). We are thus reduced to
the case in which \(X\) itself is the spectrum of a complete
noetherian local ring. But
\[
  \operatorname{Hom}_S(X,Y)=\Gamma(Z/X),
  \qquad Z=Y\times_S X,
\]
and the preceding lemma applies.

\SGAStatement{\SGAInlineRefTarget{sga3:xv:proposition:1-6}Proposition 1.6}{Let \(A,S,S_n\) be as above, and let
\(G\) and \(M\) be two \(S\)-group preschemes, with \(M\) finite over
\(S\). Then:
\begin{enumerate}
\item[\normalfont a)]\SGAInlineRefTarget{sga3:xv:proposition:1-6:item:a} The canonical map
\[
  \operatorname{Hom}_{S\text{-}\mathrm{gr}}(M,G)
  \longrightarrow
  \varprojlim_n
  \operatorname{Hom}_{S_n\text{-}\mathrm{gr}}(M_n,G_n)
\]
is bijective.
\item[\normalfont b)] If \(M\) is of multiplicative type and \(G\) is
smooth over \(S\), the canonical map
\[
  \varphi:\operatorname{Hom}_{S\text{-}\mathrm{gr}}(M,G)
  \longrightarrow
  \operatorname{Hom}_{S_0\text{-}\mathrm{gr}}(M_0,G_0)
\]
is surjective. Moreover, if
\(\varphi(u)=\varphi(u')=u_0\), then \(u\) and \(u'\) are conjugate by
an element of \(G(S)\) that reduces to the identity element of
\(G(S_0)\).
\item[\normalfont c)] If \(M\) is of multiplicative type and smooth
over \(S\), if \(G\) is flat and of finite type over \(S\), and if
\(u_0:M_0\to G_0\) is a central homomorphism, then \(u_0\) lifts
uniquely to a homomorphism
\[
  u:M\longrightarrow G.
\]
Moreover, \(u\) is central if \(G\) has connected fibres.
\end{enumerate}}
\SGARefTarget{sga3:xv:proposition:1:6}

\SGAProof{Proof}
\emph{a)} This follows from \SGARefLink{sga3:xv:corollary:1-5}{1.5}, from the fact that
\(M\times_S M\) is finite over \(S\), and from the characterization
of group morphisms as the morphisms for which the familiar diagram
\[
\begin{tikzcd}[column sep=large,row sep=large]
M\times_S M \arrow[r,"u\times u"] \arrow[d]
  & G\times_S G \arrow[d] \\
M \arrow[r,"u"'] & G
\end{tikzcd}
\]
is commutative.

\emph{b)} By Exposé \SGARefLink{sga3:ix:theorem:3:6}{IX~3.6}, one can construct a coherent system of
homomorphisms \(u_n:M_n\to G_n\) lifting a homomorphism
\(u_0:M_0\to G_0\). In view of part a), this proves the first assertion
of b).

If \(u\) and \(u'\) are two liftings of \(u_0\), then \(u_n\) and
\(u'_n\) are conjugate by an element \(g_n\in G(S_n)\) lifting the
identity element of \(G(S_0)\) (Exposé \SGARefLink{sga3:ix:theorem:3:6}{IX~3.6}). The same reference
also implies that the \(g_n\) may be chosen coherently; by \SGARefLink{sga3:xv:lemma:1-4}{1.4} they
therefore arise from a section \(g\in G(S)\). The morphisms \(u\) and
\(\operatorname{Int}(g)u'\) then agree modulo \(\mathfrak m^n\) for
every \(n\), and hence agree by \SGARefLink{sga3:xv:corollary:1-5}{1.5}.

\emph{c)} The existence and uniqueness of \(u\) follow from part a)
and \SGARefLink{sga3:xv:proposition:1:2:bis}{1.2 bis}. If \(G\) has connected fibres, \(u\) is central by
Exposé \SGARefLink{sga3:ix:theorem:5:6:item:a}{IX~5.6(a)}.

\SGASection{2}{Infinitesimal lifting of sub-tori}
\SGARefTarget{sga3:xv:section:2}

\SGASubsection{2.1}{Statement of the theorem}
\SGARefTarget{sga3:xv:subsection:2:1}

We shall give an infinitesimal lifting theorem for sub-tori of a group
prescheme that does not require smoothness hypotheses (in contrast
with Exposé \SGARefLink{sga3:ix:theorem:3:6:bis}{IX~3.6 bis}) and that also answers a very natural
question:
is it enough to be able to lift ``sufficiently many'' points of finite
order of a sub-torus \(T_0\) in order to ensure that \(T_0\) itself can
be lifted (infinitesimally, of course)?

\SGAStatement{Theorem 2.1}{Let \(S\) be a noetherian affine scheme,
let \(S_0\) be a closed subscheme of \(S\) defined by a square-zero
ideal \(J\), let \(G\) be an \(S\)-group prescheme of finite type, let
\(G_0=G\times_S S_0\), let \(T_0\) be a sub-torus of \(G_0\), and let
\(q>0\) be an integer invertible on \(S\). Suppose that, for every
integer \(n\) equal to a power of \(q\), there exists a subscheme
\(M_n\) of \(G\), flat over \(S\), such that
\[
  M_n\times_S S_0={}_nT_0.
\]
Then there exists a sub-torus \(T\) of \(G\) such that
\(T\times_S S_0=T_0\).}
\SGARefTarget{sga3:xv:theorem:2:1}

\SGARefLink{sga3:xv:theorem:2:1}{Theorem 2.1} will be useful to us through the following two
corollaries.

\SGAStatement{\SGAInlineRefTarget{sga3:xv:corollary:2-2}Corollary 2.2}{Let \(S\) be a locally noetherian
prescheme, let \(S_0\) be a closed subprescheme of \(S\) defined by a
locally nilpotent sheaf of ideals, let \(G\) be an \(S\)-group
prescheme of finite type, let \(T_0\) be a sub-torus of
\(G_0=G\times_S S_0\), and let \(q>0\) be an integer invertible on
\(S\). As \(n\) ranges over the powers of \(q\), let \((M_n)\) be a
coherent system of \(S\)-subgroup schemes of \(G\), of multiplicative
type, lifting the \({}_nT_0\). (The system \((M_n)\) is said to be
coherent if
\[
  M_m={}_m(M_n)
\]
whenever \(m\) divides \(n\).) Then there exists a unique sub-torus
\(T\) of \(G\) such that
\[
  T\times_S S_0=T_0
  \qquad\text{and}\qquad
  {}_nT=M_n
\]
for every \(n\).}
\SGARefTarget{sga3:xv:corollary:2:2}

\SGAStatement{Corollary 2.3}{Let \(G\) be an \(S\)-group prescheme,
flat and of finite presentation over \(S\), let \(S_0\) be a closed
subprescheme of \(S\) defined by a locally nilpotent sheaf of ideals
of finite type, and let \(T_0\) be a central torus of
\(G_0=G\times_S S_0\). Then there exists a unique sub-torus \(T\) of
\(G\) lifting \(T_0\). Moreover, \(T\) is central.}
\SGARefTarget{sga3:xv:corollary:2:3}

\SGAStatement{Remark 2.4}{We leave it to the reader to formulate the
analogues of statements \SGARefLink{sga3:xv:theorem:2:1}{2.1}, \SGARefLink{sga3:xv:corollary:2-2}{2.2}, and \SGARefLink{sga3:xv:corollary:2:3}{2.3} in which, instead of lifting
a sub-torus of \(G_0\), one is given a torus \(T\) over \(S\) and asks
to lift a morphism
\[
  u_0:T_0\longrightarrow G_0.
\]
One reduces to the preceding cases by considering the graph of
\(u_0\).}
\SGARefTarget{sga3:xv:remark:2:4}

Let us show how \SGARefLink{sga3:xv:corollary:2-2}{Corollaries 2.2} and \SGARefLink{sga3:xv:corollary:2:3}{2.3} follow from \SGARefLink{sga3:xv:theorem:2:1}{Theorem 2.1}.

\SGAProof{Proof of \SGARefLink{sga3:xv:corollary:2-2}{Corollary 2.2}}
The uniqueness of \(T\) follows from Exposé \SGARefLink{sga3:ix:corollary:4:8:item:b}{IX~4.8(b)} and Exposé
\SGARefLink{sga3:ix:remark:4:10}{IX~4.10}. To prove existence, we may therefore reduce to the case in
which \(S\) is affine, hence noetherian, and \(S_0\) is defined by a
square-zero ideal.

\SGAStatement{Lemma 2.5}{Let \(G\) be an \(S\)-group prescheme of
finite presentation, and let \(H\) be a subgroup scheme of \(G\), of
multiplicative type and smooth over \(S\). Then
\(\operatorname{Centr}_G(H)\) is representable by an \(S\)-group
subprescheme of \(G\), of finite presentation.}
\SGARefTarget{sga3:xv:lemma:2:5}

When \(G\) is separated over \(S\), the lemma follows from Exposé
\SGARefLink{sga3:viii:examples:6:5:item:e:01}{VIII~6.5(e)}, without any smoothness hypothesis on \(H\). In the present
case, observe that the assertion is local for the fpqc topology. We
may therefore assume that \(H\) is diagonalizable, hence of the form
\(D_S(M)\). We may also assume that \(S\) is affine and then, by EGA
IV~8, noetherian. Since \(H\) is smooth and of finite type over \(S\),
the order \(q\) of the torsion subgroup of \(M\) is invertible on
\(S\) (Exposé \SGARefLink{sga3:viii:proposition:2:1:item:e:01}{VIII~2.1(e)}). It is then immediate (cf. Exposé \SGARefLink{sga3:ix:remark:4:10}{IX~4.10})
that the subgroups \({}_nH\), with \(n\) ranging over the powers of
\(q\), are schematically dense in \(H\) (Exposé \SGARefLink{sga3:ix:definition:4:1}{IX~4.1}). But
\({}_nH\) is a completely split covering of \(S\), that is, it is
isomorphic to a finite disjoint sum of copies of \(S\). Consequently
\[
  \operatorname{Centr}_G({}_nH)=Z_n
\]
is representable as the intersection of the centralizers in \(G\) of
the sections of \({}_nH\) over \(S\). It remains to apply the following
lemma.

\SGAStatement{\SGAInlineRefTarget{sga3:xv:lemma:2-5-bis}Lemma 2.5 bis}{Let \(S\) be a noetherian prescheme, let
\(G\) be an \(S\)-group prescheme of finite type, let \(H\) be a
subgroup of \(G\) of multiplicative type, and let \((H_i)\) be an
increasing filtered family of subgroups of \(G\) of multiplicative
type. Suppose that
\[
  Z_i=\operatorname{Centr}_G(H_i)
\]
is representable by a group subprescheme of \(G\). Then the family
\((Z_i)\) is stationary. If, in addition, the \(H_i\) are schematically
dense in \(H\), then
\[
  Z_i=\operatorname{Centr}_G(H)
\]
for all sufficiently large \(i\).}
\SGARefTarget{sga3:xv:lemma:2:5:bis}

To see that the family \((Z_i)\) is stationary, it is enough to prove
that the family of underlying sets \(\operatorname{ens}(Z_i)\) is
stationary. Indeed, the stationary value will then be a closed subset
of an open subset \(U\) of \(G\). After replacing \(G\) by \(U\), we
are reduced to studying a decreasing filtered family of closed
subpreschemes of a noetherian prescheme. An elementary constructibility
argument reduces us to the case in which \(S\) is integral. We then
have to show that \(\operatorname{ens}(Z_i)\) is stationary over some
nonempty open subset of \(S\). The generic fibre of \(G\) is separated
(Exposé VI\textsubscript{A}~0.3); after restricting \(S\), we may
therefore assume that \(G\) is separated over \(S\) (EGA IV~8). Then
\(Z_i\) is closed in \(G\) (Exposé VIII~6.5(e)).

To prove the last assertion of the lemma, let \(Z\) denote the
stationary value of the family \((Z_i)\). It is clear that
\(\operatorname{Centr}_G(H)\) is a subfunctor of \(Z\); we show that
\(Z\) centralizes \(H\). Let \(E\) be the subprescheme of
\(H\times_S Z\) that is the kernel of the pair of morphisms
\[
\begin{aligned}
  H\times_S Z&\rightrightarrows G,\\
  (h,c)&\longmapsto c,\\
  (h,c)&\longmapsto hch^{-1}.
\end{aligned}
\]
The prescheme \(E\) contains \(H_i\times_S Z\) for every \(i\). The
\(H_i\) are flat over \(S\), so EGA IV~11.10.9 shows that, for every
point \(s\) of \(S\), the \((H_i)_s\) are schematically dense in
\(H_s\), and \(H_i\times_S Z\) is schematically dense in
\(H\times_S Z\). Since \(G_s\) is separated, \(E_s\) is closed in
\((H\times_S Z)_s\), and is therefore equal to it. Thus \(E\) is
closed in \(H\times_S Z\), and hence equals \(H\times_S Z\). This says
that \(Z\) centralizes \(H\).

Returning to the proof of \SGARefLink{sga3:xv:corollary:2-2}{2.2}, \SGARefLink{sga3:xv:lemma:2:5}{Lemma 2.5} shows that
\[
  Z_n=\operatorname{Centr}_G(M_n)
\]
is representable, and \SGARefLink{sga3:xv:lemma:2-5-bis}{Lemma 2.5 bis} shows that the decreasing family
of subpreschemes \((Z_n)\) is stationary; let \(Z\) be its stationary
value. The group \(Z\) contains \(T_0\) and all the \(M_n\). After
replacing \(G\) by \(Z\), we may therefore suppose that the \(M_n\)
are central.

The hypotheses of \SGARefLink{sga3:xv:theorem:2:1}{Theorem 2.1} are now satisfied, so there exists a
sub-torus \(T\) of \(G\) lifting \(T_0\). The groups \({}_nT\) and
\(M_n\) are two liftings of \({}_nT_0\); they are therefore conjugate
(Exposé \SGARefLink{sga3:ix:theorem:3:2:bis}{IX~3.2 bis}), and consequently coincide because \(M_n\) is
central. This torus \(T\) has the required properties.

\SGAProof{Proof of \SGARefLink{sga3:xv:corollary:2:3}{Corollary 2.3}}
The uniqueness of \(T\) follows from Exposé \SGARefLink{sga3:ix:theorem:5:1:bis}{IX~5.1 bis}, and its
centrality follows from Exposé \SGARefLink{sga3:ix:theorem:5:6:item:b}{IX~5.6(b)}. By the usual procedure, this
observation reduces us first to the case in which \(S\) is affine
(and hence \(S_0\) is defined by a nilpotent ideal of finite type),
and then to the noetherian case. After restricting \(S\) if necessary,
we may suppose that there exists an integer \(q\) invertible on \(S\).
\SGARefLink{sga3:xv:corollary:2:3}{Corollary 2.3} then follows from \SGARefLink{sga3:xv:corollary:2-2}{2.2} and \SGARefLink{sga3:xv:proposition:1:2}{1.2}.

\SGAStatement{Remark 2.6}{\SGARefLink{sga3:xv:corollary:2:3}{Corollary 2.3} remains true if the torus
\(T_0\) is replaced by a smooth central subgroup of multiplicative
type of \(G_0\).}
\SGARefTarget{sga3:xv:remark:2:6}

\SGASubsection{2.2}{Proof of \SGARefLink{sga3:xv:theorem:2:1}{Theorem 2.1}}
\SGARefTarget{sga3:xv:subsection:2:2}

\medskip\noindent\textbf{a) Reduction to the case
\(\boldsymbol{T_0=\mathbb G_{m,S_0}}\).}

By \SGARefLink{sga3:xv:proposition:1:1}{1.1}, we may assume that \(M_n\) is a subgroup of multiplicative
type. Using Exposé \SGARefLink{sga3:ix:theorem:3:2:bis}{IX~3.2 bis}, we may suppose that the family
\((M_n)\) is coherent (\SGARefLink{sga3:xv:corollary:2-2}{2.2}). The centralizer \(Z_n\) of \(M_n\) in
\(G\) is representable (\SGARefLink{sga3:xv:lemma:2:5}{2.5}), and the filtered family \((Z_n)\) is
stationary (\SGARefLink{sga3:xv:lemma:2-5-bis}{2.5 bis}). After replacing \(G\) by \(Z_n\) for sufficiently
large \(n\), we may therefore suppose that \(T_0\) and the \(M_n\) are
central. The uniqueness of a lifting of \(T_0\) is then guaranteed
by Exposé \SGARefLink{sga3:ix:theorem:5:1:bis}{IX~5.1 bis}.

Proceeding as in the proof of \SGARefLink{sga3:xv:proposition:1:1}{1.1}, we may suppose that there exist an
\(S\)-torus \(T\) and an \(S_0\)-isomorphism
\[
  u_0:T\times_S S_0\xrightarrow{\sim}T_0.
\]
Lifting \(u_0\) is equivalent to lifting \(T_0\). By uniqueness, it is
enough to prove the existence of a lifting of \(u_0\) after an affine
faithfully flat extension \(S'\to S\) of finite type (fpqc descent).
This allows us to suppose that
\[
  T=\mathbb G_{m,S}^{\,r}
\]
(Exposé \SGARefLink{sga3:x:corollary:4-5}{X~4.5}). If the restriction of \(u_0\) to each factor
\(\mathbb G_m\) lifts to an \(S\)-morphism---necessarily central---then
these liftings immediately give a lifting of \(u_0\). In short, we may
suppose that \(T_0=\mathbb G_{m,S_0}\).

\medskip\noindent\textbf{b) Definition of the obstruction to the
existence of a lifting of \(\boldsymbol{T_0}\).}

By \SGARefLink{sga3:xv:proposition:1:1}{1.1}, to prove \SGARefLink{sga3:xv:theorem:2:1}{Theorem 2.1} it is enough to find a subscheme of \(G\),
flat over \(S\), lifting \(T_0\). We shall see that the obstruction to
the existence of such a lifting can be defined as an element of a
certain \(\operatorname{Ext}^1\) group of sheaves of modules.

Let \(U\) be an open subset of \(G\) in which \(T_0\) is closed, and
again denote by \(U\) (respectively, \(U_0\)) the open subscheme of
\(G\) (respectively, \(G_0\)) having this open subset as its underlying
space. The sheaf \(\mathcal O_{T_0}\), regarded as a sheaf on \(U\), is
therefore a quotient of \(\mathcal O_{U_0}\). Let
\[
  h:\mathcal O_{U_0}\longrightarrow\mathcal O_{T_0}
\]
be the canonical epimorphism.

\SGAStatement{\SGAInlineRefTarget{sga3:xv:lemma:2-7}Lemma 2.7}{The canonical map
\[
  h'=\operatorname{id}_J\otimes h:
  J\otimes_{S_0}\mathcal O_{U_0}
  \longrightarrow
  J\otimes_{S_0}\mathcal O_{T_0}
\]
factors, necessarily uniquely, as \(i\circ j_U\), where \(j_U\) is
the canonical epimorphism
\[
  J\otimes_{S_0}\mathcal O_{U_0}
  \longrightarrow
  J\mathcal O_U\simeq J\mathcal O_{U_0}.
\]}
\SGARefTarget{sga3:xv:lemma:2:7}

We must show that \(h'(K)=0\), where \(K\) is the kernel of \(j_U\).
For every integer \(n\) equal to a power of \(q\), there is an
epimorphism
\[
  h_n:\mathcal O_{T_0}\longrightarrow\mathcal O_{{}_nT_0}.
\]
Since \({}_nT_0\) lifts to a scheme \(M_n\) flat over \(S\), the
canonical morphism
\[
  j_n:J\otimes_{S_0}\mathcal O_{{}_nT_0}
  \xrightarrow{\sim}J\mathcal O_{M_n}
\]
is an isomorphism. The commutative diagram
\[
\begin{tikzcd}[column sep=large,row sep=large]
K \arrow[r]
  & J\otimes_{S_0}\mathcal O_{U_0}
      \arrow[r,"j_U"] \arrow[d,"h'"]
  & J\mathcal O_U \arrow[r,hook] \arrow[dl,"i"]
  & \mathcal O_U \\
  & J\otimes_{S_0}\mathcal O_{T_0}
      \arrow[d,"h'_n"] & & \\
  & J\otimes_{S_0}\mathcal O_{{}_nT_0}
      \arrow[r,"j_n"',"\sim"] & J\mathcal O_{M_n}
      \arrow[r,hook] & \mathcal O_{M_n}
\end{tikzcd}
\]
shows that \(h'(K)\) is contained in \(\ker h'_n\) for every \(n\).
It is therefore contained in \(\bigcap_n\ker h'_n\), and it remains
to prove that this intersection is zero. Now
\[
  \mathcal O_{T_0}=\mathcal O_{S_0}[\mathbb Z],
  \qquad
  \mathcal O_{{}_nT_0}
     =\mathcal O_{S_0}[\mathbb Z/n\mathbb Z].
\]
Let
\[
  a=\sum_{m\in\mathbb Z}a_m\otimes m
  \in J\otimes_{S_0}\mathcal O_{T_0}.
\]
The \(a_m\) are sections of \(J\), almost all zero. Choose \(n\) large
enough that the indices \(m\) for which \(a_m\ne0\) have distinct
images in \(\mathbb Z/n\mathbb Z\). If \(a\in\ker h'_n\), then
necessarily \(a=0\). Thus \(\bigcap_n\ker h'_n=0\), proving \SGARefLink{sga3:xv:lemma:2-7}{2.7}.

Let \(K_0\) be the kernel of \(h\). Consider the diagram
\[
\begin{tikzcd}[column sep=large,row sep=large]
  &&& 0 \arrow[d] \\
  & J\mathcal O_U \arrow[d,"\sim"] && K_0 \arrow[d] \\
0 \arrow[r]
  & J\mathcal O_{U_0} \arrow[r] \arrow[d,"i"']
  & \mathcal O_U \arrow[r] & \mathcal O_{U_0}
      \arrow[r] \arrow[d,"h"'] & 0 \\
  & J\otimes_{S_0}\mathcal O_{T_0}
  && \mathcal O_{T_0} \arrow[d] \\
  &&& 0 .
\end{tikzcd}
\]
The sheaf \(\mathcal O_U\) defines an element
\[
  \Phi\in
  \operatorname{Ext}^{1}_{\mathcal O_U}
  (\mathcal O_{U_0},J\mathcal O_{U_0}).
\]
Let \(\Psi\) be the element of
\[
  \operatorname{Ext}^{1}_{\mathcal O_U}
  (K_0,J\otimes_{S_0}\mathcal O_{T_0})
\]
obtained from \(\Phi\), by the bifunctoriality of
\(\operatorname{Ext}^{1}\), through the morphisms
\[
  K_0\longrightarrow\mathcal O_{U_0},
  \qquad
  i:J\mathcal O_{U_0}
    \longrightarrow J\otimes_{S_0}\mathcal O_{T_0}.
\]
It follows from Exposé \SGARefLink{sga3:iii:lemma:4-1}{III~4.1} and the infinitesimal criterion of
flatness (cf. Exposé \SGARefLink{sga3:iii:proposition:4-3}{III~4.3}) that a subscheme of \(U\), flat over
\(S\) and lifting \(T_0\), exists if and only if \(\Psi=0\).

Since \(T_0\) is affine, Exposé \SGARefLink{sga3:iii:proposition:4-5}{III~4.5} and \SGARefLink{sga3:xv:lemma:4:6}{4.6} show that it is enough
to construct, locally on \(U\), a subscheme flat over \(S\) lifting
\(T_0\). It therefore suffices to show that the image of \(\Psi\) in
the sheaf
\[
  \mathscr E=
  \mathcal Ext^{1}_{\mathcal O_U}
  (K_0,J\otimes_{S_0}\mathcal O_{T_0})
\]
is zero. We again denote this element of
\(\Gamma(U,\mathscr E)\) by \(\Psi\).

\medskip\noindent\textbf{c) Reduction to the case in which \(S\) is
local artinian with algebraically closed residue field.}

Since \(K_0\) and \(J\otimes_{S_0}\mathcal O_{T_0}\) are coherent
sheaves, so is \(\mathscr E\); moreover, \(\mathscr E\) is supported
on \(T_0\), because the same is true of
\(J\otimes_{S_0}\mathcal O_{T_0}\). To show that the section \(\Psi\)
of \(\mathscr E\) is zero, it is enough to show that, for every point
\(u\) of \(T_0\), its image in the fibre of \(\mathscr E\) at \(u\)
is zero. Formation of the \(\mathcal Ext^i\) of coherent sheaves
commutes with flat extensions of the base, so we are reduced to proving the
existence of a lifting of
\[
  T_0\cap\operatorname{Spec}\mathcal O_u
\]
flat over \(S\). Let \(s\) be the image of \(u\) in \(S\). We may
replace \(S\) by \(\operatorname{Spec}\mathcal O_s\) and \(G\) by
\(G\times_S\operatorname{Spec}\mathcal O_s\).

After another faithfully flat extension if necessary, we may suppose
that \(\mathcal O_s\) has an algebraically closed residue field
(EGA \(0_{\mathrm{III}}\), 10.3.1). Let \(\mathfrak m\) be the maximal
ideal of \(\mathcal O_s\), and set
\[
  S_n=\operatorname{Spec}(\mathcal O_s/\mathfrak m^n).
\]
Suppose that, for every \(n>0\), the obstruction to lifting
\[
  T_0\times_S S_n=(T_0)_n
\]
to a subscheme of \(G_n=G\times_S S_n\), flat over \(S_n\), has been
shown to vanish, and let \(T_n\) be the unique flat lifting of
\((T_0)_n\) over \(S_n\) that is a sub-torus of \(G_n\). If
\(n>n'\), it is clear that
\[
  T_n\times_{S_n}S_{n'}=T_{n'}.
\]
If \(u\) is a point of \(T_0\) lying over \(s\), the lemma below,
applied to the coherent system of liftings
\[
  T_n\cap\operatorname{Spec}\mathcal O_u
\]
of
\((T_0)_n\cap\operatorname{Spec}\mathcal O_u\), shows that
\(T_0\cap\operatorname{Spec}\mathcal O_u\) indeed has a lifting flat
over \(\mathcal O_s\). We are therefore reduced to proving that
\(\Psi=0\) when \(S=S_n\) is the spectrum of an artinian local ring
with algebraically closed residue field, and to proving the following
lemma.

\SGAStatement{\SGAInlineRefTarget{sga3:xv:lemma:2-8}Lemma 2.8}{Let \(A\to B\) be a local homomorphism of
noetherian local rings, let \(\mathfrak m\) be the maximal ideal of
\(A\), let \(J\) be a square-zero ideal of \(A\), let \(M\) be a
finite-type \(B\)-module, and set
\[
  A_0=A/J,\qquad B_0=B/JB,\qquad M_0=M/JM.
\]
Let \(N_0\) be a quotient \(B_0\)-module of \(M_0\), flat over \(A_0\).
For every \(n>0\), let \(A_n,A_{0,n}\), and so forth denote the objects
obtained by the base extension
\[
  A\longrightarrow A/\mathfrak m^n=A_n,
\]
and let \(\overline{J}_n\) be the image of \(J\) in \(A_n\). For every \(n>0\),
let \(N_n\) be a quotient \(B_n\)-module of \(M_n\), flat over \(A_n\),
lifting \(N_{0,n}\), and suppose that, for \(n\ge n'\), \(N_{n'}\) is
obtained from \(N_n\) by the base extension \(A_n\to A_{n'}\). Then
there exists a quotient \(B\)-module \(N\) of \(M\), flat over \(A\),
lifting \(N_0\).}
\SGARefTarget{sga3:xv:lemma:2:8}

\SGAProof{Proof of \SGARefLink{sga3:xv:lemma:2-8}{2.8}}
Let \(P_0\) be the kernel of the epimorphism \(M_0\to N_0\). For every
\(n>0\), there is a commutative diagram with exact horizontal rows:
\[
\tag{$\ast_n$}\label{sga3:xv:lemma:2-8:equation:ast-n}
\begin{tikzcd}[column sep=2em,row sep=large]
&&&& 0 \arrow[d] & \\
&&&& P_{0,n} \arrow[d] & \\
& 0 \arrow[r]
  & \overline{J}_nM_n \arrow[r] \arrow[dd]
  & M_n \arrow[r] \arrow[dd]
  & M_{0,n} \arrow[r] \arrow[dd] & 0 \\
\overline{J}_n\otimes_{A_{0,n}}M_{0,n}
      \arrow[urr] \arrow[drr] &&&&& \\
& 0 \arrow[r]
  & \overline{J}_n\otimes_{A_{0,n}}N_{0,n} \arrow[r]
  & N_n \arrow[r] \arrow[d]
  & N_{0,n} \arrow[r] \arrow[d] & 0 \\
&&& 0 & 0 &
\end{tikzcd}
\]
Moreover, by hypothesis, the diagram \((\ast)_{n+1}\) reduces modulo
\(\mathfrak m^n\) to \((\ast)_n\).

The Artin--Rees lemma (Bourbaki, \emph{Algèbre commutative}, Chapter
III, §3, Corollary 1) shows that the filtrations on \(JM\),
\(J\otimes_{A_0}M_0\), and \(J\otimes_{A_0}N_0\), defined by the
kernels of the respective morphisms
\[
\begin{gathered}
  JM\longrightarrow \overline{J}_nM_n,\\
  J\otimes_{A_0}M_0
    \longrightarrow \overline{J}_n\otimes_{A_{0,n}}M_{0,n},\\
  J\otimes_{A_0}N_0
    \longrightarrow \overline{J}_n\otimes_{A_{0,n}}N_{0,n},
\end{gathered}
\]
are \(\mathfrak mB\)-good. Passing to the inverse limit in the
diagrams \((\ast)_n\), and writing \(\widehat Q\) for the separated
completion of a \(B\)-module \(Q\) for the \(\mathfrak mB\)-adic
topology, we obtain a commutative diagram in which the two horizontal
rows remain exact:
\[
\tag{$\widehat\ast$}\label{sga3:xv:lemma:2-8:equation:widehat-ast}
\begin{tikzcd}[column sep=large,row sep=large]
&&&& 0 \arrow[d] & \\
&&&& \widehat P_0 \arrow[d] & \\
& 0 \arrow[r]
  & \widehat{JM} \arrow[r] \arrow[dd]
  & \widehat M \arrow[r] \arrow[dd]
  & \widehat M_0 \arrow[r] \arrow[dd] & 0 \\
\widehat{J\otimes_{A_0}M_0}
      \arrow[urr] \arrow[drr] &&&&& \\
& 0 \arrow[r]
  & \widehat{J\otimes_{A_0}N_0} \arrow[r]
  & \varprojlim_n N_n \arrow[r] \arrow[d]
  & \widehat N_0 \arrow[r] \arrow[d] & 0 \\
&&& 0 & 0 &
\end{tikzcd}
\]
Now \((B,\mathfrak mB)\) is a Zariski ring and
\(J\otimes_{A_0}N_0\) is of finite type over \(B\), hence separated
for the \(\mathfrak mB\)-adic topology. The diagram
\((\widehat\ast)\) then shows that the morphism
\[
  J\otimes_{A_0}M_0\longrightarrow J\otimes_{A_0}N_0
\]
induced by the epimorphism \(M_0\to N_0\) factors through \(JM\):
\[
\begin{tikzcd}[column sep=large,row sep=large]
J\otimes_{A_0}M_0 \arrow[rr,"\mathrm{can.}"] \arrow[dr]
  && JM \arrow[dl] \\
& J\otimes_{A_0}N_0 &
\end{tikzcd}
\]

Under these conditions, Exposé \SGARefLink{sga3:iii:lemma:4-1}{III~4.1} and \SGARefLink{sga3:iii:proposition:4-3}{4.3} show that a quotient
\(B\)-module \(N\) of \(M\), flat over \(A\) and lifting \(N_0\),
exists if and only if a certain element
\[
  g\in E=\operatorname{Ext}^{1}_B
  (P_0,J\otimes_{A_0}N_0)
\]
vanishes. More precisely, the exact sequence
\[
  0\longrightarrow JM\longrightarrow M
  \longrightarrow M_0\longrightarrow0
\]
defines an element \(f\in F\), where
\[
  F=\operatorname{Ext}^{1}_B(M_0,JM),
\]
and \(g\) is the image of \(f\) under the natural morphism \(F\to E\)
arising by bifunctoriality from
\[
  P_0\longrightarrow M_0,
  \qquad
  JM\longrightarrow J\otimes_{A_0}N_0.
\]
The diagram \((\widehat\ast)\) and Exposé \SGARefLink{sga3:iii:lemma:4-1}{III~4.1} show that the image
of \(g\) in
\[
  \widehat E\simeq
  \operatorname{Ext}^{1}_{\widehat B}
  (\widehat P_0,\widehat{J\otimes_{A_0}N_0})
\]
is zero. But \(E\) is of finite type over \(B\), and hence separated
for the \(\mathfrak mB\)-adic topology. Consequently \(g=0\), which
completes the proof of \SGARefLink{sga3:xv:lemma:2-8}{2.8}.

\medskip\noindent\textbf{d) Study of \(\boldsymbol{\mathscr E}\).}

We now suppose that \(S\) is the spectrum of an artinian local ring
\(A\) with algebraically closed residue field \(k\). Let \(A_0\) be
the ring of \(S_0\), and let \(\mathfrak m_0\) be the maximal ideal
of \(A_0\). Since \(S_0\) is artinian, \(T_0\) is closed in \(G\)
(Exposé VI\textsubscript{B}~\SGARefLink{cor:sga3-VIB-1.4.2}{1.4.2}). We may therefore take \(U=G\), so
that
\[
  \mathscr E=
  \mathcal Ext^{1}_{\mathcal O_G}
  (K_0,J\otimes_{S_0}\mathcal O_{T_0}).
\]
Let \(B_0\) be the ring of the affine \(S_0\)-scheme
\(T_0=\mathbb G_{m,S_0}\), and let \(\overline B_0\) be the ring of its special
fibre
\[
  T_0\times_{S_0}\operatorname{Spec}(k)=\mathbb G_{m,k}.
\]
The sheaf \(\mathscr E\) is a coherent
\(\mathcal O_{T_0}\)-module, hence is defined by a finite-type
\(B_0\)-module, which we denote by \(E\).

Consider the graded module associated with \(E\) and the
\(\mathfrak m_0B_0\)-adic filtration:
\[
  E_r=\mathfrak m_0^rE/\mathfrak m_0^{r+1}E.
\]
Each \(E_r\) is a finite-type \(\overline B_0\)-module, and \(E_r=0\) for all
sufficiently large \(r\), since \(S_0\) is artinian.

Let \(g\) be a section of \(G\) over \(S\) whose restriction to
\(S_0\) is a section \(g_0\) of \(T_0\). Left translation by \(g\)
defines an ``automorphism of the situation'' from the point of view
of the obstruction problem. In particular, \(g\) induces an
\(S_0\)-automorphism of the sheaf \(\mathscr E\) that fixes the
obstruction \(\Psi\). More precisely, \(g\) defines a semilinear
automorphism of the \(B_0\)-module \(E\), relative to the
\(A_0\)-automorphism of \(B_0\) defined by translation by \(g_0\) in
\(T_0\). Reducing modulo \(\mathfrak m_0^{r+1}\), \(g\) therefore
defines a semilinear automorphism of \(E_r\), relative to the
\(k\)-automorphism of \(\overline B_0\) defined by translation by
\[
  g_0\times_{S_0}\operatorname{Spec}(k)
\]
in \(T_0\times_{S_0}\operatorname{Spec}(k)\).

\SGAStatement{Lemma 2.9}{For every integer \(r\ge0\), \(E_r\) is a
locally free \(\overline B_0\)-module.}
\SGARefTarget{sga3:xv:lemma:2:9}

Let \(x\) be a point of \(T_0\), let \(\kappa(x)\) be its residue
field, and let
\[
  (E_r)_x=E_r\otimes_{\overline B_0}\kappa(x)
\]
be the ``fibre'' of \(E_r\) at \(x\). Denote by \(\ell(x)\) the rank
of \((E_r)_x\) over \(\kappa(x)\), and by \(\ell\) the maximum of
\(\ell(x)\) as \(x\) ranges over the points of \(T_0\). Let \(L\) be
the largest closed subscheme of
\[
  \operatorname{Spec}\overline B_0=\mathbb G_{m,k}
\]
over which \(E_r\) is locally free of rank \(\ell\) (TDTE IV, \SGARefLink{sga3:xv:lemma:3:6}{Lemma
3.6}). Let \(\beta\in L(k)\), which exists because \(L\) is of finite
type over the algebraically closed field \(k\), and let
\(\alpha\in\mathbb G_{m,k}(k)\) be a point whose order is a power
\(n\) of \(q\). The point \(\alpha\) is \(k\)-rational. Since, by
hypothesis, \({}_nT_0\) lifts to a subscheme \(M_n\) étale over \(S\),
there is a section \(a\) of \(G\) over \(S\) lifting \(\alpha\) and
whose restriction to \(S_0\) is a section of \(T_0\).

The preceding observations show that the fibres of \(E_r\) at the
points \(\beta\) and \(\alpha\beta\) of \(\mathbb G_{m,k}(k)\) are
\(k\)-isomorphic. The points whose order is a power of \(q\) are dense
in \(\mathbb G_{m,k}\), as are their translates by \(\beta\). Since
\(L\) is closed in \(\mathbb G_{m,k}\) and \(\mathbb G_{m,k}\) is
reduced, \(L=\mathbb G_{m,k}\), and \(E_r\) is locally free of rank
\(\ell\) over \(\overline B_0\).

\medskip\noindent\textbf{e) Completion of the proof of \SGARefLink{sga3:xv:theorem:2:1}{Theorem 2.1}.}

Let \(K_0(n)\) be the sheaf of ideals of \(\mathcal O_{G_0}\)
defining the closed subscheme \({}_nT_0\). The sheaf
\(K_0\) is
therefore a subsheaf of \(K_0(n)\). For brevity, set
\[
  R=J\otimes_{S_0}\mathcal O_{T_0},
  \qquad
  R(n)=J\otimes_{S_0}\mathcal O_{{}_nT_0}.
\]
The sheaf \(R(n)\) is a quotient of \(R\), and we have the diagram
\[
\begin{tikzcd}[column sep=large,row sep=large]
&& K_0 \arrow[d] \\
&& K_0(n) \arrow[d] \\
0 \arrow[r]
  & J\mathcal O_G \arrow[r] \arrow[d]
  & \mathcal O_G \arrow[r]
  & \mathcal O_{G_0} \arrow[r] & 0 \\
  & R \arrow[d] \\
  & R(n)
\end{tikzcd}
\]
Using the bifunctoriality of
\(\operatorname{Ext}^{1}_{\mathcal O_G}\), we obtain the commutative
diagram
\[
\begin{tikzcd}[column sep=huge,row sep=large]
\mathcal{E}xt^{1}_{\mathcal O_G}
  (\mathcal O_{G_0},J\mathcal O_G)
  \arrow[dr] \\
&
\mathcal{E}xt^{1}_{\mathcal O_G}(K_0(n),R)
  \arrow[r] \arrow[d]
&
\mathcal{E}xt^{1}_{\mathcal O_G}(K_0,R)=\mathscr E
  \arrow[d] \\
&
\mathcal{E}xt^{1}_{\mathcal O_G}(K_0(n),R(n))
  \arrow[r]
&
\mathcal{E}xt^{1}_{\mathcal O_G}(K_0,R(n)).
\end{tikzcd}
\]
Again let \(\Phi\) denote the element of
\(\operatorname{Ext}^{1}_{\mathcal O_G}
(\mathcal O_{G_0},J\mathcal O_G)\) defined by the exact sequence
\[
  0\longrightarrow J\mathcal O_G
  \longrightarrow\mathcal O_G
  \longrightarrow\mathcal O_{G_0}
  \longrightarrow0,
\]
so that \(\Psi\) is the image of \(\Phi\) in \(\mathscr E\). Since
\({}_nT_0\) lifts to a subscheme \(M_n\) of \(G\), flat over \(S\),
the image of \(\Phi\) in
\[
  \operatorname{Ext}^{1}_{\mathcal O_G}(K_0(n),R(n))
\]
is zero (Exposé \SGARefLink{sga3:iii:lemma:4-1}{III~4.1}). A fortiori, the image of \(\Phi\) in
\(\operatorname{Ext}^{1}_{\mathcal O_G}(K_0,R(n))\), which is also
the image of \(\Psi\), is zero.

\SGAStatement{Lemma 2.10}{For every integer \(n\), the canonical
morphism
\[
  \operatorname{Ext}^{1}_{\mathcal O_G}(K_0,R)
    \otimes_{B_0}\mathcal O_{{}_nT_0}
  \longrightarrow
  \operatorname{Ext}^{1}_{\mathcal O_G}(K_0,R(n))
\]
is injective.}
\SGARefTarget{sga3:xv:lemma:2:10}

Indeed, the affine scheme \(T_0=\mathbb G_{m,S_0}\) has ring
\[
  B_0=A_0[T,T^{-1}].
\]
The subscheme \({}_nT_0\) is defined by the vanishing of
\[
  h(n)=T^n-1.
\]
This section is regular (EGA \(0_{\mathrm{IV}}\), 15.2.2) and remains
regular after every base change \(S'_0\to S_0\). Hence there is an
exact sequence of sheaves
\[
  0\longrightarrow\mathcal O_{T_0}
  \xrightarrow{\,h(n)\,}\mathcal O_{T_0}
  \longrightarrow\mathcal O_{{}_nT_0}
  \longrightarrow0.
\]
Since \({}_nT_0\) is flat over \(S\), tensoring over
\(\mathcal O_{S_0}\) with \(J\) gives the exact sequence
\[
  0\longrightarrow R\xrightarrow{\,h(n)\,}R
  \longrightarrow R(n)\longrightarrow0,
\]
and then the exact sequence of Ext groups
\[
  \cdots\longrightarrow
  \operatorname{Ext}^{1}_{\mathcal O_G}(K_0,R)
  \xrightarrow{\,h(n)\,}
  \operatorname{Ext}^{1}_{\mathcal O_G}(K_0,R)
  \longrightarrow
  \operatorname{Ext}^{1}_{\mathcal O_G}(K_0,R(n)),
\]
which proves the lemma.

The foregoing shows that, for every integer \(n\) equal to a power of
\(q\), the image of \(\Psi\) in
\(\mathscr E/h(n)\mathscr E\) is zero. To prove that \(\Psi=0\), it
is enough to show that
\[
  \Psi\in\mathfrak m_0^r\mathscr E
  \quad\Longrightarrow\quad
  \Psi\in\mathfrak m_0^{r+1}\mathscr E.
\]
Let \(\overline\Psi\) be the image of \(\Psi\) in \(E_r\). There is
an element \(\Psi(n)\in\mathscr E\) such that
\[
  \Psi=\Psi(n)h(n)
  \qquad(n\text{ a power of }q).
\]
We observed that the image \(\overline{h(n)}\) of \(h(n)\) in \(\overline B_0\)
is again \(\overline B_0\)-regular. Since \(E_s\) is locally free over \(\overline B_0\)
for every \(s\) (\SGARefLink{sga3:xv:lemma:2:9}{2.9}), multiplication by \(\overline{h(n)}\) on
\(E_s\) is injective. It follows immediately that
\(\Psi(n)\in\mathfrak m_0^r\mathscr E\). Let
\(\overline{\Psi(n)}\) be its image in \(E_r\); then
\[
  \overline\Psi=\overline{h(n)}\,\overline{\Psi(n)}
  \qquad(n\text{ a power of }q).
\]
Thus the set \(F\) of points of \(\mathbb G_{m,k}\) at which
\(\overline\Psi\) vanishes contains the dense set of points whose
order is a power of \(q\). Moreover, \(F\) is closed because \(E_r\)
is locally free over \(\overline B_0\), and \(\mathbb G_{m,k}\) is reduced.
Therefore \(\overline\Psi=0\).

This completes the proof of \SGARefLink{sga3:xv:theorem:2:1}{Theorem 2.1}.

